\chapter{Discussion and Conclusions}
\label{ch:conclusions}

\section{The instrument}

This thesis set out to measure whether a quantized spiking neural network
reasons like its full-precision parent, and found that the measurement could not
be interpreted until the instrument was calibrated.

A four-block convolutional spiking network classifies four families of transient
noise in LIGO data at $0.9833 \pm 0.0014$ validation accuracy over three seeds,
under constraints --- shift-implementable membrane decay, at most eight time
steps, 112-pixel input --- that keep the configuration deployable by
construction rather than by subsequent rounding. Each of those constraints was
shown to cost nothing measurable in accuracy, and the resolution choice improved
the health of the deepest layer by 12.7 percentage points of dead fraction.

The Spike Activation Map was implemented, its summation range corrected, and its
free parameter fixed on the temporal axis it governs rather than on
time-aggregated maps that cannot see it. Its faithfulness was tested
quantitatively and found to hold on two classes of four, with a geometric
explanation for the other two.

The calibration established three levels: the divergence produced by weight noise
as large as an entire 8-bit quantization step (PCC $0.975$), by a thirtyfold
change in the explainer (PCC $0.960$), and by initialization alone (PCC $0.153$).
The third is far larger than expected, and it carries three consequences.

\paragraph{Spatial metrics cannot compare independently trained models.} Two
networks that agree to within one validation sample in accuracy produce
explanations correlating at $0.153$. The mechanism is representational
permutation rather than any defect of the metric, and the disagreement is uniform
across samples rather than driven by outliers.

\paragraph{They can compare a model with its own parent.} Under a perturbation 85
times smaller than the 2-bit quantization step, PCC remains above $0.975$. The
two results answer different questions, and the distinction must be made
explicitly whenever such a number is reported. In this thesis it decided the
design of the grid: every quantized model is fine-tuned from its own parent.

\paragraph{Published figures are not interpretable without these floors.} Values
around $0.85$ are presented in the literature as evidence of preserved
explanations. For this architecture that value falls between the numerical floor
and the independent-solutions level, so it cannot be read without knowing which
comparison produced it. Establishing the floors requires no training runs beyond
a multi-seed baseline that is good practice regardless.

A fourth observation cuts across all of them: the class sensitivity ordering is
identical under three unrelated interventions. Explanation sensitivity is a
property of the class morphology, not of the perturbation applied, and
Extremely\_Loud --- saturated and spatially extensive --- registers no change
under any measure tested.

The preparation of the membrane arm added a result of its own: the network makes
no use of negative membrane potentials. Clipping the hidden state to $[0, 2]$,
which removes up to 47 thresholds of negative range, costs one validation sample
and leaves the dynamics unchanged, so the membrane can be stored on an unsigned
grid at one bit less than a symmetric one would need.

\section{Quantization and explanation fidelity}

\paragraph{Accuracy is not a sufficient criterion for spiking networks either.}
The finding of Rabbi et al.~\cite{rabbi2026} on convolutional networks carries
over, in a sharper form. Two configurations of this network classify equally
well, a 4-bit weight model and a 2-bit membrane model, and their explanations
differ by a factor of almost two in agreement with the original. Anyone selecting
a configuration on accuracy alone could have chosen the second and deployed a
model that reaches the right labels while showing a different part of the
spectrogram. In the detector-characterization setting of
Chapter~\ref{ch:introduction}, that is the scenario that matters: the label is
correct, the diagnostic cue is not the one the original model would have given.

\paragraph{The divergence is real but bounded.} The changed explanations are
still faithful: deleting the pixels they highlight degrades the model's
confidence almost as much as deleting the parent's highlights does, and the
quantized model can be explained with its parent's maps as well as with its own.
The quantized network has not stopped using spectrogram evidence; among
several redundant sets of evidence it now highlights a different one. A
faithfulness test alone would have declared the explanation intact; an overlap
test alone would have suggested it was lost. Both are needed to say what
happened.

\paragraph{The hypothesis was wrong in its form and right in its mechanism.}
The prediction that the membrane would move the temporal axis of the explanation
before the weights moved its spatial axis is not supported: both arms first
register at 4 bits, and at 2 bits the membrane changes where the explanation
points more than when. What the membrane does change first is the temporal
dynamics of the network itself: at 2 bits, with a single representable level
between rest and threshold, the first layer fires with perfectly regular
intervals. The temporal centre of mass, built to capture the time axis of the
explanation, is blind to this, because it measures where the spike mass falls,
not how it is spaced. A temporal explanation metric for spiking networks needs
to see regularity as well as timing, and the present instrument does not.

\paragraph{The comparison with Smith et al.\ inverts.} Smith et
al.~\cite{smith2026} report that the membrane needs at least four bits while
learned weights survive two. The weight result agrees with theirs: uniform ternary weights fail, where they
needed a learned quantizer to keep two bits. The membrane result does not: with
a grid fitted to the measured range, a 2-bit membrane keeps accuracy at parity,
where they lose $4$--$8\,\%$.

\paragraph{Weak events are not where the explanation fails.} The expectation that
quantization would disturb spike decisions mostly in weak events, whose
potentials sit near threshold, finds no consistent support: the gradient with
event strength changes sign across seeds. The only consistent effect is
reversed: the most saturated Extremely\_Loud events lose their explanation most,
because their evidence is spread over the whole image and the highlighted region
is least constrained.

\paragraph{Efficiency and explainability do not trade off here.} The
configuration that changes the explanation at constant accuracy is the
2-bit membrane, and it is the one that saves no energy: with weights in full
precision the per-operation cost stays the same while firing rates rise. The
deployable configuration, joint 4-bit, keeps accuracy at parity and the
explanation at the numerical floor, at about one sixtieth of the estimated
energy and one eighth of the memory. The divergence between the three curves of
Figure~\ref{fig:headline} exists, but it occurs at a point no one would choose
for efficiency.

\paragraph{Seed variance is the dominant source of variation.} Three
full-precision networks of equal accuracy differ by a factor of 2.6 in the
firing rate of their first layer, by more than half in estimated energy, and in
which classes SAM can explain faithfully at all. The 2-bit membrane IoU ranges
from $0.26$ to $0.62$ across seeds. Every comparison in this thesis was made
within a seed, and a single-seed study of the same question could have reached
any conclusion from ``the explanation is preserved'' to ``it is lost''.

\section{Limitations}

\paragraph{Generalization of the thresholds.} The three calibration levels are
specific to this architecture, this dataset and this explainer. The
\emph{mechanism} behind the collapse, representational permutation between
independently trained networks, is general and well documented; the
\emph{magnitude} is not, and a different architecture could place the floors
elsewhere. What generalizes is the claim that the floors must be measured, not
the values measured here.

\paragraph{Labels are model output.} Accuracy is agreement with the Gravity Spy
convolutional classifier rather than with physical ground truth. A model that
reproduced that classifier's systematic errors would score well here. The
per-class confidences are available and would permit stratifying accuracy by the
reference classifier's own certainty, which was not done.

\paragraph{The partition is random with respect to time.} If glitches of the same
class are correlated within observing periods, a random test set overestimates
generalization to new periods. A block split by observing run would be the more
conservative validation; the timestamps were dropped during feature selection for
the reasons given in Chapter~\ref{ch:methodology} and would have to be recovered.

\paragraph{$T$ sits at the imposed ceiling.} The search was forbidden from
exploring more than eight time steps, so it is not known whether it would have
preferred a longer integration window. An unconstrained search selected seven,
which bounds the likely cost but does not eliminate the question.

\paragraph{The deletion metric fails on two classes.} The geometric explanation
offered in Section~\ref{sec:faithfulness} --- that deletion-based tests penalize
spatially extended evidence --- is an interpretation, not a demonstration. An
insertion curve, which rewards rather than penalizes distributed evidence, would
test it directly.

\paragraph{The explanation is class-agnostic.} SAM does not use labels, so the
deletion metric measures faithfulness to the model's \emph{predicted} class. This
is the correct reading of the metric and matches its original formulation, but it
means misclassified examples are informative and should be inspected separately
rather than pooled.

\paragraph{The temporal metric is partly kernel arithmetic.} The $\gamma$
analysis showed that the centre of mass of the SAM energy is dominated by a
minority of near-continuously firing neurons. Between models compared at the same
$\gamma$ this component cancels, but it means that the CoM is most sensitive to
changes in that minority and may under-report changes elsewhere.

\paragraph{Quantization is simulated.} Fake quantization reproduces the values an
integer implementation would hold, not its arithmetic: the membrane is quantized
as a stored state after each spike decision, while on integer hardware the
comparison with threshold would itself be performed on integers. The energy
figures are analytical, at a 45\,nm reference node, and below eight bits they
rest on an extrapolation of the per-operation costs.

\paragraph{Sample sizes of the explanation analyses.} Divergence is measured on
twenty validation images per class, faithfulness on ten, stratification on
fifteen per class and tercile, and the membrane range on a single checkpoint.
These are sufficient for the effects reported here, which are large, but limit
the resolution with which small differences between configurations can be
ranked; the faithfulness analysis covers only the 2-bit membrane arm.

\paragraph{One grid, one fitting procedure.} Each bit-width was fitted once per
seed, with one fine-tuning schedule and one membrane range. The collapse of the
2-bit weight arm and the survival of the 2-bit membrane are properties of these
quantizers; a different weight quantizer (learned scales, a wider grid for the
readout) could move the weight arm's breaking point.

\paragraph{The control does not isolate the fine-tuning.} Fine-tuning without a
quantizer never improved the parent's validation accuracy, so the selected
control checkpoint is the parent itself. The control therefore shows that the
pipeline is inert, but not what the extra epochs alone would do to a model that
does change; the quantized checkpoints combine both effects.

\paragraph{Checkpoints are selected on the validation partition} on which they
are then analyzed. With selection by accuracy and analysis of explanations, the
dependence is weak, but the test-set evaluation is the only fully independent
estimate.

\paragraph{The saving factor rests on an assumption.} Costing the analog first
convolution at the weight bit-width gives a joint 4-bit saving of 62; costing it
in floating point gives 3.5. The operation counts, the memory and the ranking of
configurations do not depend on it.

\section{Future work}

\paragraph{Deployment.} Hardware deployment was deliberately left out of scope.
The natural continuation is to export the selected integer configuration to a
neuromorphic platform or a microcontroller, and to measure on it the latency,
energy and memory that are estimated analytically here.

\paragraph{Normalization.} The baseline contains no batch normalization, because
quantizing its parameters is a problem distinct from quantizing weights and
potentials. Deeper spiking networks need it, and both approximating it with shifts
and removing it altogether through progressive training are available; repeating
the study on a normalized architecture would test whether the conclusions survive.

\paragraph{The temporal axis of the input.} Under direct encoding every temporal
structure is generated by the network. Feeding the spectrogram column by column,
so that time in the input corresponds to time in the signal, would give the
temporal explanation a physical meaning and make the temporal metrics measure
something a detector scientist can read.

\paragraph{Wider benchmark.} Extending the four classes to the full Gravity Spy
taxonomy, with a split by observing run, would make the classification results
comparable with the published benchmark.

\section{Closing remark}

Compressed models are increasingly deployed where their explanations are acted
upon. This thesis asked whether a quantized spiking network still decides for
the same reasons as the network it came from, and found that the question has
two prerequisites: knowing what the comparison metrics report when nothing has
changed, and comparing every model with its own parent. With both in place, the
answer for this network is that it mostly does --- down to four bits, in every
arm --- and that the one configuration that does not is invisible to accuracy,
partly invisible to correlation and to faithfulness, and visible to a metric
that asks which pixels a reader is shown. That is the configuration a
deployment chosen on accuracy could have picked, and the reason the question was
worth asking.
